% TODO make hw stylesheet

\documentclass[10pt]{article}
\usepackage[letterpaper,top=1in,left=1in,right=1in]{geometry}
\usepackage[dvipsnames,svgnames]{xcolor}
\usepackage[shortlabels]{enumitem}
\usepackage[framemethod=TikZ]{mdframed}
\usepackage{amsmath,amssymb,amsthm}
\usepackage[colorlinks]{hyperref}
\usepackage{multicol}
\usepackage{tabularx}
\usepackage{bbm}
\usepackage{tikz-cd}

\hypersetup{
    colorlinks=true,
    linkcolor=blue,
    filecolor=magenta,
    urlcolor=magenta,
    pdftitle={MAT 535 HW}
}

\usepackage{mathtools}
\usepackage{thmtools}
\usepackage[textsize=footnotesize]{todonotes}
\usepackage{titling}
\usepackage{cancel}

\reversemarginpar
\mdfdefinestyle{mdbluebox}{roundcorner=10pt,innerbottommargin=9pt,
linecolor=blue,backgroundcolor=TealBlue!5,}
\declaretheoremstyle[headfont=\sffamily\bfseries\color{MidnightBlue},
mdframed={style=mdbluebox},]{thmbluebox}

\mdfdefinestyle{mdbloobox}{frametitlefont=\bfseries,innerbottommargin=8pt,
nobreak=true,backgroundcolor=TealBlue!5,linecolor=blue,}
\declaretheoremstyle[headfont=\bfseries\color{MidnightBlue},
mdframed={style=mdbloobox},headpunct={\\[3pt]},postheadspace=0pt,]{thmbloobox}

\mdfdefinestyle{mdredbox}{frametitlefont=\bfseries,innerbottommargin=8pt,
nobreak=true,backgroundcolor=Salmon!5,linecolor=RawSienna,}
\declaretheoremstyle[headfont=\bfseries\color{RawSienna},
mdframed={style=mdredbox},headpunct={\\[3pt]},postheadspace=0pt,]{thmredbox}

\mdfdefinestyle{mdgreenbox}{linecolor=ForestGreen,backgroundcolor=ForestGreen!5,
linewidth=2pt,rightline=false,leftline=true,topline=false,bottomline=false,}
\declaretheoremstyle[headfont=\bfseries\sffamily\color{ForestGreen!70!black},
mdframed={style=mdgreenbox},headpunct={ --- },]{thmgreenbox}

\mdfdefinestyle{mdorangebox}{linecolor=RedOrange,backgroundcolor=RedOrange!5,
linewidth=2pt,rightline=false,leftline=true,topline=false,bottomline=false,}
\declaretheoremstyle[headfont=\bfseries\sffamily\color{RedOrange!70!black},
mdframed={style=mdorangebox},headpunct={ --- },]{thmorangebox}

\mdfdefinestyle{mdblackbox}{linecolor=black,backgroundcolor=RedViolet!5!gray!5,
linewidth=3pt,rightline=false,leftline=true,topline=false,bottomline=false,}
\declaretheoremstyle[mdframed={style=mdblackbox}]{thmblackbox}

\declaretheoremstyle[spaceabove=6pt,spacebelow=6pt]{boldhead}
\declaretheoremstyle[spaceabove=6pt,spacebelow=6pt,headfont=\normalfont\itshape]{italichead}

\declaretheorem[style=boldhead,name=Problem]{problem}
\declaretheorem[style=boldhead,name=Problem,numbered=no]{problem*}
\declaretheorem[style=boldhead,name=Setup]{setup} % for config geo
\declaretheorem[style=boldhead,name=Example,sibling=problem]{example}
\declaretheorem[style=boldhead,name=$\bigstar$ Problem,sibling=problem]{niceprob}
\declaretheorem[style=boldhead,name=Theorem,sibling=problem]{theorem}
\declaretheorem[style=boldhead,name=Corollary,sibling=theorem]{corollary}
\declaretheorem[style=boldhead,name=Proposition,sibling=theorem]{prop}
\declaretheorem[style=boldhead,name=Lemma,numberwithin=problem]{lemma}
\declaretheorem[style=boldhead,name=Theorem,numbered=no]{theorem*}
\declaretheorem[style=boldhead,name=Lemma,numbered=no]{lemma*}
\declaretheorem[style=boldhead,name=Claim,numberwithin=problem]{claim}
\declaretheorem[style=boldhead,name=Claim,numbered=no]{claim*}
\declaretheorem[style=boldhead,name=Remark,numberwithin=problem]{remark}
\declaretheorem[style=boldhead,name=Remark,numbered=no]{remark*}
\declaretheorem[style=boldhead,name=Definition,sibling=theorem]{definition}
\declaretheorem[style=boldhead,name=Definition,numbered=no]{definition*}

\providecommand{\ol}{\overline}
\providecommand{\tensor}{\otimes}
\renewcommand{\hom}{\operatorname{Hom}}
\providecommand{\tor}{\operatorname{Tor}}
\providecommand{\ceil}[1]{\left\lceil #1 \right\rceil}
\providecommand{\floor}[1]{\left\lfloor #1 \right\rfloor}
\newcommand{\dd}{\,\mathrm{d}}
\providecommand{\ul}{\underline}
\providecommand{\eps}{\varepsilon}
\providecommand{\half}{\frac{1}{2}}
\providecommand{\CC}{\mathbb C}
\providecommand{\FF}{\mathbb F}
\providecommand{\II}{\mathbb I}
\providecommand{\NN}{\mathbb N}
\providecommand{\QQ}{\mathbb Q}
\providecommand{\RR}{\mathbb R}
\providecommand{\ZZ}{\mathbb Z}
\providecommand{\ind}{\mathbbm 1}
\providecommand{\dg}{^\circ}
\providecommand{\ii}{\item}
\providecommand{\setdiff}{\smallsetminus}
\providecommand{\alert}[1]{{\sffamily\textbf{\textcolor{blue}{#1}}}}
\providecommand{\inv}{^{-1}}
\providecommand{\mb}[1]{\mathbf{#1}}
\newcommand{\what}{\widehat}

\newenvironment{soln}{\begin{proof}[Solution]}{\end{proof}}

\providecommand{\pow}{\operatorname{Pow}}
\providecommand{\vocab}[1]{{\textbf{\textcolor{ForestGreen}{#1}}}}
\providecommand{\lcm}{\operatorname{lcm}}
\providecommand{\abs}[1]{\left\lvert #1\right\rvert}
\providecommand{\smabs}[1]{\lvert #1\rvert}
\providecommand{\innerprod}[1]{\langle\!\langle #1\rangle\!\rangle}
\providecommand{\norm}[1]{\left\| #1\right\|}
\providecommand{\smnorm}[1]{\| #1\|}
\providecommand{\gr}{\operatorname{gr}}

\usepackage{mathrsfs}

% place title at top right
\setlength{\droptitle}{-8ex}
\pretitle{\begin{flushright}\Large\bfseries}
\posttitle{\par\end{flushright}}
\preauthor{\begin{flushright}}
\postauthor{\end{flushright}}
\predate{\begin{flushright}}
\postdate{\end{flushright}}

\title{MAT 535 HW04}
\author{\large Aditya Pahuja}
\date{\large Due 03/05}
\begin{document}
\maketitle

\begin{problem}
    [\S13.1\#2]
    Show that $x^3 - 2x - 2$ is irreducible over $\QQ$
    and let $\theta$ be a root.
    Compute $(1+\theta)(1+\theta+\theta^2)$
    and $\frac{1+\theta}{1+\theta+\theta^2}$ in $\QQ(\theta)$.
\end{problem}

\begin{soln}
    If $x^3 - 2x - 2$ factors over $\QQ$, it must have
    a linear factor, implying that it has a rational root.
    The rational root theorem says that a rational root
    of the given polynomial must be one of $\pm 1$ or $\pm 2$,
    but we can check manually that none of these are roots
    ($\pm 1$ fail because the value of the polynomial
    would be odd, while $\pm 2$ fail because
    the value of the polynomial would be $2\bmod 4$).

    We can compute that
    \begin{equation*}
	(1 + \theta)(1 + \theta  +\theta^2)
	= 1 + 2\theta + 2\theta^2 + \theta^3
	= \boxed{3 + 4\theta + 2\theta^2}
    \end{equation*}
    and
    \begin{equation*}
	\frac{1+\theta}{1+\theta+\theta^2}
	= (1+\theta)\cdot \frac13(5 + \theta - 2\theta^2)
	= \frac13(5 + 6\theta - \theta^2 - 2\theta^3)
	= \boxed{\frac13 + \frac23\theta - \frac13\theta^2}
    \end{equation*}
\end{soln}

\begin{problem}
    [\S13.2\#4]
    Determine the degree over $\QQ$ of
    $2+\sqrt 3$ and of $1 + \sqrt[3]2 + \sqrt[3]4$.
\end{problem}

\begin{soln}
    The degree of $2+\sqrt3$ is $\boxed 2$;
    clearly the degree is greater than $1$ because $\sqrt3$,
    hence $2+\sqrt3$, is not rational,
    and the degree is at most $2$ because
    $2 + \sqrt3$ is a root of $x^2 - 4x + 1$.

    The degree of $\alpha = 1+\sqrt[3]2 + \sqrt[3]4$ is $\boxed 3$.
    Obviously $\alpha\in\QQ(\sqrt[3]2)$, hence
    $\QQ(\alpha)\subseteq\QQ(\sqrt[3]2)$,
    and we can compute
    \begin{equation*}
	\alpha(\sqrt[3]2 - 1) = (\sqrt[3]2)^3 - 1 = 1
	\iff \frac 1\alpha + 1 = \sqrt[3]2,
    \end{equation*}
    so the reverse inclusion $\QQ(\sqrt[3]2)\subseteq\QQ(\alpha)$ also holds.
    Thus $\sqrt[3]2$ has the same degree as $\alpha$, namely $3$.
\end{soln}

\begin{problem}
    [\S13.2\#8]
    Let $F$ be a field with characteristic not equal to $2$.
    Let $D_1$ and $D_2$ be elements of $F$,
    neither of which is a square in $F$.
    Prove that $F(\sqrt{D_1}, \sqrt{D_2})$ is of degree $4$ over $F$
    if $D_1D_2$ is not a square in $F$ and is of degree $2$ over $F$ otherwise.
\end{problem}

\begin{soln}
    We show that $F(\sqrt{D_1}) = F(\sqrt{D_2})$
    if and only if $D_1D_2$ is a square in $F$.

    First, if $D_1D_2$ is a square in $F$, then
    there exists $a\in F$ such that
    $\sqrt{D_1}\sqrt{D_2} = a$, so
    $\sqrt{D_1} = a/\sqrt{D_2}\in F(\sqrt{D_2})$
    and $\sqrt{D_2} = a/\sqrt{D_1}\in F(\sqrt{D_1})$
    implies that $F(\sqrt{D_1})$ and $F(\sqrt{D_2})$ contain each other.

    Conversely, if $F(\sqrt{D_1}) = F(\sqrt{D_2})$,
    then $\sqrt{D_1} = a_1 + a_2\sqrt{D_2}$
    for $a_1,a_2\in F$. Squaring,
    \begin{equation*}
	D_1 = (a_1^2 + D_2a_2^2) + 2a_1a_2\sqrt{D_2}
    \end{equation*}
    implying that $2a_1a_2\sqrt{D_2}\in F$,
    hence $2a_1a_2 = 0$. Since $F$ does not have characteristic $2$,
    either $a_1 = 0$, in which case $\sqrt{D_1} = a_2\sqrt{D_2}$
    implies that $D_1D_2 = (a_2D_2)^2$ is a square,
    or $a_2 = 0$, in which case $\sqrt{D_1} = a_1\in F$,
    which contradicts the fact that $D_1$ is not a square;
    thus it must be the case that $D_1D_2$ is a square in $F$.

    Thus if $D_1D_2$ is a square then
    $F(\sqrt{D_1},\sqrt{D_2}) = F(\sqrt{D_1})$
    has degree $2$. If $D_1D_2$ is not a square then
    \begin{enumerate}[(1)]
	\ii $F(\sqrt{D_1},\sqrt{D_2})$ has degree
	divisible by the least common multiple of
	$[F(\sqrt{D_1}) : F]$ and $[F(\sqrt{D_2}) : F]$, namely $2$.
	\ii $F(\sqrt{D_1},\sqrt{D_2})$ has degree greater than $2$,
	since $F(\sqrt{D_1})$ is a proper subfield of degree $2$.
	\ii $F(\sqrt{D_1},\sqrt{D_2})$ has degree at most
	$[F(\sqrt{D_1}) : F][F(\sqrt{D_2}) : F] = 4$.
    \end{enumerate}
    These three pieces of numerical information
    imply that $[F(\sqrt{D_1},\sqrt{D_2}) : F] = 4$.
\end{soln}

\begin{problem}
    [\S13.2\#16]
    Let $K/F$ be an algebraic extension
    and let $R$ be a ring contained in $K$ containing $F$.
    Show that $R$ is a subfield of $K$ containing $F$.
\end{problem}

\begin{soln}
    For any $\alpha\in R$,
    if $p(x) = c_0 + c_1x + \cdots + c_nx^n \in F[x]$
    is the minimal polynomial of $\alpha$ over $F$
    (which exists because $\alpha$ is algebraic over $F$),
    then we can write
    \begin{equation*}
	\alpha\inv = -c_0\inv(c_1 + c_2\alpha + \cdots + c_n\alpha^{n-1}).
    \end{equation*}
    Since $R$ contains $F$ and $\alpha$ it must contain
    the coefficients $-c_0\inv c_k\in F$
    and the powers of $\alpha$, so $R$ contains $\alpha\inv$
    for every $\alpha\in R$, hence $R$ is a field
    contained in $K$ and containing $F$.
\end{soln}

\begin{problem}
    [\S13.2\#20]
    Show that if the matrix of the linear transformation
    ``multplication by $\alpha$'' considered in the previous exercise
    is $A$ then $\alpha$ is a root of the characteristic polynomial of $A$.
    This gives an effective procedure for determining
    an equation of degree $n$ satisfied by an element $\alpha$
    in an extension of $F$ of degree $n$.
    Use this procedure to obtain the monic polynomial
    of degree $3$ satisfied by $\sqrt[3]2$
    and by $1+\sqrt[3]2 + \sqrt[3]4$.
\end{problem}

\begin{soln}
    Suppose that $[F(\alpha) : F] = d$ and $[K : F(\alpha)] = n/d$.
    The extension $F(\alpha)/F$ has $F$-basis
    $1$, $\alpha$, \dots, $\alpha^{d-1}$,
    and $K/F(\alpha)$ has some $F(\alpha)$-basis
    $\beta_1$, \dots, $\beta_{n/d}$, so that
    the elements $\alpha^j \beta_k\in K$
    for $0\le j < d$ and $1\le k \le n/d$ form an $F$-basis of $K$.
    Multiplication by $\alpha$ sends
    $\alpha^j\beta_k$ to $\alpha^{j+1}\beta_k$
    for $0\le j < d-1$ and sends $\alpha^d\beta_k$ to
    $(-c_0-c_1\alpha-\cdots-c_{d-1}\alpha^{d-1})\beta_k$
    where $p(x) = c_0 + c_1x + \cdots + c_dx^d\in F[x]$
    is the minimal polynomial of $\alpha$ over $x$,
    so the matrix of $\alpha$ on the invariant subspace
    generated by $\beta_k$, $\alpha\beta_k$, \dots, $\alpha^{d-1}\beta_k$ is
    \begin{equation*}
        \begin{pmatrix}
	    0 & 0 & 0 & \cdots & -c_0\\
	    1 & 0 & 0 & \cdots & -c_1\\
	    0 & 1 & 0 & \cdots & -c_1\\
	    \vdots & \vdots & \vdots & \ddots & \vdots\\
	    0 & 0 & 0 & \cdots & -c_{d-1}\\
        \end{pmatrix},
    \end{equation*}
    the companion matrix $M_{p(x)}$ of $p(x)$.
    Then the matrix of $A$ with respect to our basis is
    \begin{equation*}
	\begin{pmatrix}
	    M_{p(x)} & 0 & \cdots & 0 \\
	    0 & M_{p(x)} & \cdots & 0 \\
	    \vdots & \vdots & \ddots & \vdots \\
	    0 & 0 & \cdots & M_{p(x)}
	\end{pmatrix}
    \end{equation*}
    where the $0$s are $d\times d$ blocks of zeroes
    and $M_{p(x)}$ appears $n/d$ times.
    The characteristic polynomial of this matrix is
    $p(x)^{n/d}$, which of course has $\alpha$ as a root.

    Observe that $\alpha = \sqrt[3]2$ and $\beta = 1 + \sqrt[3]2 + \sqrt[3]4$
    are both in $F(\sqrt[3]2)$, which has basis
    $\{1,\sqrt[3]2,\sqrt[3]4\}$. With respect to this basis,
    multplication by $\alpha$ has matrix
    \begin{equation*}
        \begin{pmatrix}
	    0 & 0 & 2\\
	    1 & 0 & 0\\
	    0 & 1 & 0\\
        \end{pmatrix}
    \end{equation*}
    which has characteristic polynomial $\boxed{x^3 - 2}$;
    multplication by $\beta$ has matrix
    \begin{equation*}
        \begin{pmatrix}
	    1 & 2 & 2\\
	    1 & 1 & 2\\
	    1 & 1 & 1
        \end{pmatrix}
    \end{equation*}
    which has characteristic polynomial
    $\boxed{x^3 - 3x^2 - 3x - 1}$.
\end{soln}










\end{document}
